\chapter{Methodology} \label{ch:methodology}

The experiments are structured around a simple logic: a bare baseline is defined first,
and one technique is added at a time. The Qwen 2.5 model family is used, tested across
four sizes from 0.5B to 7B. Five techniques are covered: constrained decoding,
quantization, few-shot prompting, RAG, and LoRA fine-tuning. All configurations are
evaluated on the Berkeley Function Calling Leaderboard.

\section{Experimental Design}
The experimental design is based on controlled comparisons. A baseline configuration
is established first, and each technique is added one at a time. This isolates the
effect of each technique on tool-calling accuracy.

\subsection{Cumulative Evaluation Ladder}

Five configurations are evaluated, each building on the previous one. Config B is
the baseline: the model is prompted without constrained decoding at full precision.
In Config CD, constrained decoding is added, forcing the output to conform to a
valid JSON schema. In Config CD+Q, AWQ INT4 quantization is applied to reduce memory footprint
while preserving accuracy as much as possible. In Config CD+Q+FT, LoRA
fine-tuning is added to the quantized model. Two parallel configurations are also
evaluated: Config PE uses few-shot prompting, and Config CD+Q+RAG adds
retrieval-augmented generation to Config CD+Q.

\begin{figure}[ht]
\centering
\begin{tikzpicture}[
    box/.style={
        rectangle, draw, rounded corners=4pt,
        minimum width=3.0cm, minimum height=0.9cm,
        align=center, font=\small
    },
    main/.style={box, fill=blue!8},
    side/.style={box, fill=orange!8},
    para/.style={box, fill=gray!10, draw=gray!60},
    arr/.style={-{Stealth}, thick},
    darr/.style={-{Stealth}, thick, dashed, gray!70}
]

\node[main] (B)     {Config B};
\node[main, right=1.6cm of B]   (CD)    {Config CD};
\node[main, right=1.6cm of CD]  (CDQ)   {Config CD+Q};
\node[main, right=1.6cm of CDQ] (CDQFT) {Config CD+Q+FT};

\node[side, below=1.2cm of CDQ]  (CDQRAG) {Config CD+Q+RAG};
\node[para, below=1.2cm of B]    (PE)     {Config PE};

\draw[arr] (B)   -- node[above, font=\footnotesize]{$+$CD}  (CD);
\draw[arr] (CD)  -- node[above, font=\footnotesize]{$+$Q}   (CDQ);
\draw[arr] (CDQ) -- node[above, font=\footnotesize]{$+$FT}  (CDQFT);
\draw[arr] (CDQ) -- node[right, font=\footnotesize]{$+$RAG} (CDQRAG);
\draw[darr](B)   -- node[right, font=\footnotesize]{parallel} (PE);

\end{tikzpicture}
\caption{Cumulative evaluation ladder. Each configuration in the main chain (top row)
extends the previous one by adding a single technique. Config CD+Q+RAG and Config PE
are evaluated as parallel tracks.}
\label{fig:eval-ladder}
\end{figure}

\subsection{Hardware and Infrastructure}

\section{Evaluation Framework}
\subsection{Berkeley Function Calling Leaderboard}
\subsection{Metrics}

\section{Models}
\subsection{Qwen 2.5}

\section{Baseline Definition}

\section{Constrained Decoding}

\section{Quantization}

\section{Prompt Engineering and Few-Shot}

\section{Retrieval-Augmented Generation}

\section{LoRA Fine-Tuning}
